\subsection*{从角动量定义到可直接使用的判断}

角动量的参考点不能省略。同一质点在同一时刻具有唯一的动量 $\vect p$，但相对于不同参考点，
位置矢量 $\vect r$ 不同，因而角动量也不同。把参考点由 $O$ 改为 $O'$，若
$\overrightarrow{OO'}=\vect a$，则
\begin{equation}
 \vect L_{O'}=(\vect r-\vect a)\times\vect p
 =\vect L_O-\vect a\times\vect p.
\end{equation}
所以写角动量守恒以前，应先固定参考点，并检查外力对该点的总力矩是否为零。碰撞题常把
参考点选在支点或瞬时接触点，就是为了使未知的巨大冲力通过参考点，从而不贡献冲量矩。

角动量的大小也可写成
\begin{equation}
 L=p\,d,
\end{equation}
其中 $d$ 是参考点到动量作用线的垂直距离。这一形式与力矩 $M=Fd$ 完全平行：动量描述
转动状态，力矩描述角动量的改变。若力在有限时间内作用，对角动量定理积分可得
\begin{equation}
 \Delta\vect L=\int_{t_1}^{t_2}\vect M\,\dd t,
\end{equation}
右侧称为冲量矩。短时碰撞中真正应被比较的是外力的冲量矩，而不是外力瞬时值；即使某个
外力很大，只要作用时间很短或力臂很小，它对碰撞前后角动量关系的影响仍可能忽略。

外积的正负方向必须和所选坐标系一致。在平面问题中，通常规定纸面向外为正，则
\begin{equation}
 L_z=xp_y-yp_x,
 \qquad
 M_z=xF_y-yF_x.
\end{equation}
此时逆时针趋势为正、顺时针趋势为负。这个符号规则比逐次使用右手螺旋更方便，但不能在同一
道题中途改变。若所有运动都在一个平面内，角动量和力矩虽然仍是矢量，却只需保留垂直平面的
一个有符号分量。

\subsection*{有心力运动的径向与横向分解}

把速度分解成径向分量 $\dot r$ 和横向分量 $r\dot\varphi$，质点相对于力心的角动量为
\begin{equation}
 L=mr^2\dot\varphi.
\end{equation}
于是单位时间扫过的面积
\begin{equation}
 \frac{\dd A}{\dd t}=\frac12r^2\dot\varphi=\frac{L}{2m}
\end{equation}
保持不变。这就是开普勒面积定律背后的动力学原因：并非只有平方反比引力才有面积速度守恒，
任何始终指向固定力心的力都满足这一结论。

有心保守力问题还可写成一维径向运动。总能量为
\begin{equation}
 E=\frac12m\dot r^2+\frac{L^2}{2mr^2}+U(r)
 =\frac12m\dot r^2+U_{\rm eff}(r),
\end{equation}
其中
\begin{equation}
 U_{\rm eff}(r)=U(r)+\frac{L^2}{2mr^2}
\end{equation}
称为有效势。第二项并不是新的真实相互作用，而是角动量守恒对径向运动造成的运动学限制。
回归点满足 $\dot r=0$，也就是 $E=U_{\rm eff}(r)$；圆轨道还要求
$\dd U_{\rm eff}/\dd r=0$。若极小值存在，极小值附近的小幅径向振动是稳定的；若只有极大值，
对应圆轨道对扰动不稳定。

这套方法也说明为什么回归点处不能只写机械能守恒。能量方程给出速度大小，角动量方程给出
横向运动与半径的联系；两条守恒律同时使用，才能从两个回归点的半径求出两个速度，或反过来
判断轨道能否到达某一区域。

\subsection*{质点系、刚体与转动惯量}

质点系角动量的质心分解可以理解成“整体公转加内部自转”。分解式中的交叉项消失，依赖的是
质心定义
\begin{equation}
 \sum_i m_i\vect r_i'=\vect0,
 \qquad
 \sum_i m_i\vect v_i'=\vect0.
\end{equation}
它并不要求系统是刚体，也不要求各质点之间的距离不变。因此粒子团、爆炸碎片和可形变物体
同样可以使用这一分解；只有把自转角动量进一步写成 $I\vect\omega$ 时，才需要刚体、定轴或
主轴转动等附加条件。

转动惯量不仅取决于总质量和物体形状，还取决于转轴。若已知刚体绕通过质心且与目标轴平行的
轴转动惯量为 $I_C$，两轴距离为 $d$，则平行轴定理给出
\begin{equation}
 \boxed{I=I_C+Md^2}.
\end{equation}
证明时令任一质元到质心轴的横向位置为 $\vect\rho_i$，到新轴的横向位置为
$\vect\rho_i+\vect d$。展开平方后，中间项因 $\sum m_i\vect\rho_i=0$ 消失。这个结果解释了
为什么同一根杆绕端点转动比绕中心转动更“难”：转轴平移使全部质量平均离轴更远。

对于薄平面刚体，若 $x,y$ 轴位于板面内并通过同一点，$z$ 轴垂直板面，则垂直轴定理为
\begin{equation}
 I_z=I_x+I_y.
\end{equation}
它只适用于所有质元确实处在同一平面内的薄板，不能直接用于具有厚度的三维物体。

\subsection*{角动量守恒与陀螺进动的边界}

系统总角动量守恒并不要求每个部分的角动量分别守恒。人收拢双臂时，肌肉做功改变了转动动能，
但若外力矩可忽略，总角动量仍守恒：$I$ 变小，$\omega$ 增大。由
\begin{equation}
 E_k=\frac{L^2}{2I}
\end{equation}
可见在 $L$ 不变时，$I$ 变小反而使转动动能增大，增加的能量来自人体内力做功。角动量守恒
与机械能守恒是两条独立的条件，不能互相替代。

稳定进动公式 $\omega_{\rm p}=Mgr/(I_0\omega_0)$ 是高速自转近似。它把总角动量近似看成
沿陀螺对称轴，并忽略章动以及进动本身对角动量的贡献。若自转不够快，轴倾角会明显振荡，
不能再把重力矩简单等同于水平角动量方向的匀速转动率。使用该式时还应确认支点近似固定、
接触处无明显滑动，并把 $r$ 取为支点到质心的距离，而不是陀螺的几何半径。
